\section{Higgs bundles on foliated Kähler manifolds}
\subsection{Harmonic Higgs bundles}
We keep the notation of the previous sections. In particular, $G_{\mathbb R}$ is a simple real algebraic Hermitian group, $K_{\mathbb R}$ its maximal compact subgroup, $M=G_{\mathbb R}/K_{\mathbb R}$ the associated irreducible Hermitian symmetric space of the noncompact type, $G$ and $K$ are the complexifications of $G_{\mathbb R}$ and $K_{\mathbb R}$, and $\mathfrak g=\mathfrak k\oplus\mathfrak m$ and $\mathfrak g_{\mathbb R}=\mathfrak k_{\mathbb R}\oplus\mathfrak m_{\mathbb R}$ are the associated Cartan decompositions of the Lie algebras of $G$ and $G_{\mathbb R}$.

Let now $(Y,\omega_Y)$ be a compact Kähler manifold, $\Gamma=\pi_1(Y)$ its fundamental group, and $\rho:\Gamma\to G_{\mathbb R}$ be a \emph{reductive representation} (group homomorphism) of $\Gamma$ into $G_{\mathbb R}$. The assumption that $\rho$ is reductive means that the real Zariski closure of $\rho(\Gamma)$ in $G_{\mathbb R}$ is a reductive group.

In this case, by \cite{jep93Cor88}, there exists a $\rho$-equivariant harmonic map $f$ from the universal cover $\widetilde Y$ of $Y$ to the symmetric space $M=G_{\mathbb R}/K_{\mathbb R}$ associated with $G_{\mathbb R}$. The fact that $Y$ is Kähler and the nonpositivity of the complexified sectional curvature of $M$ imply by a Bochner formula due to \cite{jep93Sam78,jep93Siu80} that the map $f$ is pluriharmonic (i.e., its restriction to 1-dimensional complex submanifolds of $Y$ is still harmonic), and that the image of its (complexified) differential at every point $y\in Y$ is an abelian subalgebra of $T^{\mathbb C}_{f(y)}M$ identified with $\mathfrak m$.

By the work of Hitchin and Simpson, this gives a \emph{harmonic $G_{\mathbb R}$-Higgs principal bundle} $(P_K,\theta)$ on $Y$. We will now briefly describe the construction and the properties of such a Higgs bundle. Details and proofs can be found in the original papers \cite{jep93Hit87,jep93Hit92,jep93Sim88,jep93Sim92}.

There is a flat principal bundle $P_{G_{\mathbb R}}\to Y$ of group $G_{\mathbb R}$ associated with the representation $\rho$. The $\rho$-equivariant map $f:\widetilde Y\to G_{\mathbb R}/K_{\mathbb R}$ defines a reduction of its structure group to $K_{\mathbb R}$, i.e., a principal $K_{\mathbb R}$ bundle $P_{K_{\mathbb R}}\subset P_{G_{\mathbb R}}$. The differential of $f$ can be seen as a 1-form with values in $P_{K_{\mathbb R}}(\mathfrak m_{\mathbb R}):=(P_{K_{\mathbb R}}\times\mathfrak m_{\mathbb R})/K_{\mathbb R}$, the vector bundle associated with the adjoint action of $K_{\mathbb R}$ on $\mathfrak m_{\mathbb R}$.

If we enlarge the structure group of $P_{K_{\mathbb R}}$ to $K$, the pluriharmonicity of $f$ implies that the $K$-principal bundle $P_K\to Y$ is a holomorphic bundle and that the (1,0)-part $d^{1,0}f:T^{1,0}Y\to T^{\mathbb C}M$ of the complexified differential of $f$ defines a holomorphic section $\theta$ of $P_K(\mathfrak m)\otimes\Omega_Y^1$, where $P_K(\mathfrak m)$ is the holomorphic vector bundle associated with the principal bundle $P_K$ and the adjoint representation of $K$ on $\mathfrak m$. The section $\theta$ is called the \emph{Higgs field} and satisfies the integrability condition $[\theta,\theta]=0$ as a section of $P_K(\mathfrak m)\otimes\Omega_Y^2$. The pair $(P_K,\theta)$ is called a $G_{\mathbb R}$-Higgs principal bundle on $Y$, see e.g. \cite[\S2.2]{jep93BGPG06}.

If now $\jepPubBbb{E}$ is a (complex) representation of $G$ we can construct the associated holomorphic vector bundle $E:=P_K(\jepPubBbb{E})$ over $Y$. Since $\jepPubBbb{E}$ is a representation of $\mathfrak g$ and not only of $\mathfrak k$, we have a morphism $P_K(\mathfrak m)\to P_K(\operatorname{End}(\jepPubBbb{E}))=\operatorname{End}(E)$. Thus, the Higgs field $\theta$ can be seen as a holomorphic 1-form with values in the endomorphisms of $E$, i.e., a section of $\operatorname{End}(E)\otimes\Omega_Y^1$. The pair $(E,\theta)$ is called a \emph{$G_{\mathbb R}$-Higgs vector bundle} on $Y$. The harmonic map $f$, seen as a reduction of the structure group of $P_K$ to the compact subgroup $K_{\mathbb R}$, together with a $K_{\mathbb R}$-invariant metric on $\jepPubBbb{E}$, gives a Hermitian metric on $(E,\theta)$ called the \emph{harmonic metric}.

The existence of this harmonic metric and the fact that $P_K$ comes from a flat principal $G_{\mathbb R}$ bundle imply that for any representation $\jepPubBbb{E}$ of $G$, the associated Higgs vector bundle $(E,\theta)\to Y$ is Higgs polystable of degree 0, see \cite{jep93Sim88}. To explain what Higgs polystability means, we first define Higgs subsheaves of the Higgs bundle $(E,\theta)$. A coherent subsheaf $\jepPubScript{F}$ of $\jepPubScript{O}_Y(E)$ is a \emph{Higgs subsheaf} if it is invariant by the Higgs field, i.e., if it satisfies $\theta(\jepPubScript{F})\subset\jepPubScript{F}\otimes\Omega_Y^1$. The Higgs vector bundle $(E,\theta)$ is said to be \emph{Higgs stable} if for any Higgs subsheaf $\jepPubScript{F}$ of $(E,\theta)$ such that $0<\operatorname{rk}\jepPubScript{F}<\operatorname{rk}E$, we have $\mu_{\omega_Y}(\jepPubScript{F})<\mu_{\omega_Y}(E)$, where $\mu_{\omega_Y}(\jepPubScript{F})$ is the slope of $\jepPubScript{F}$, i.e., its degree (computed w.r.t. the Kähler form $\omega_Y$ of $Y$) divided by its rank. The Higgs bundle $(E,\theta)$ is \emph{Higgs polystable} if it is a direct sum of Higgs stable Higgs vector bundles of the same slope. Note that here $E$ is flat as a $C^{\infty}$ bundle, so that its degree is zero.
\begin{remark}
Since moreover we assumed that $G_{\mathbb R}$ is a Hermitian group, then as a $K$-representation we have $\mathfrak m=\mathfrak m^+\oplus\mathfrak m^-$ and the Higgs field $\theta$ on the principal bundle $P_K$ (or on any associated vector bundle $E$) has two components $\beta\in P_K(\mathfrak m^-)\otimes\Omega_Y^1$ and $\gamma\in P_K(\mathfrak m^+)\otimes\Omega_Y^1$. The vanishing of $\gamma$, resp. $\beta$, means that the harmonic map $f$ is holomorphic, resp. antiholomorphic. The component $\beta$, resp. $\gamma$, will be called the holomorphic, resp. antiholomorphic, component of the Higgs field $\theta$.
\end{remark}
\subsection{Harmonic Higgs bundles on foliated Kähler manifolds}
Assume now that the base Kähler manifold $Y$ of the harmonic $G_{\mathbb R}$-Higgs vector bundle $(E,\theta)\to Y$ of degree $0$ admits a smooth holomorphic foliation by complex curves and that this foliation $\jepPubScript{T}$ admits a transverse volume form $\Omega$. This means that $\Omega$ is a nowhere vanishing semipositive closed $(d-1,d-1)$-form ($d$ is the dimension of $Y$) such that the interior product of $\Omega$ with any vector tangent to the foliation $\jepPubScript{T}$ is zero.

Our goal in this section is to briefly explain the interplay between the Higgs bundle and the foliation (equipped with its transverse volume form). Details can be found in \cite[\S2.2]{jep93KM17}.

We first weaken the notion of Higgs subsheaves of $(E,\theta)$ to \emph{leafwise Higgs subsheaves} by asking only an invariance by the Higgs field along the leaves. More precisely we now consider the Higgs field as a section of $\operatorname{Hom}(E\otimes L,E)$, where $L$ is the holomorphic line subbundle of $T_Y$ tangent of the foliation $\jepPubScript{T}$. A leafwise Higgs subsheaf $\jepPubScript{F}$ of $(E,\theta)$ is a subsheaf of $\jepPubScript{O}_Y(E)$ such that $\theta(\jepPubScript{F}\otimes L)\subset\jepPubScript{F}$.

We define the foliated degree $\deg_{\jepPubScript{T}}\jepPubScript{F}$ and the foliated slope $\mu_{\jepPubScript{T}}(\jepPubScript{F})$ of a coherent sheaf $\jepPubScript{F}$ on $Y$ by wedging the first Chern class of $\jepPubScript{F}$ with the transverse volume form $\Omega$:
\[
\deg_{\jepPubScript{T}}\jepPubScript{F}:=\int_Y c_1(\jepPubScript{F})\wedge\Omega,\quad\text{and}\quad\mu_{\jepPubScript{T}}(\jepPubScript{F})=\frac{\deg_{\jepPubScript{T}}\jepPubScript{F}}{\operatorname{rk}\jepPubScript{F}}.
\]
The link between the notions of slope (and stability) for modules and for sheaves used in this paper will be given by Fact~4.5 below.

In order to state the needed leafwise polystability property of the Higgs bundle $(E,\theta)$ w.r.t. the foliated degree, we first recall a few definitions.

An open subset of $Y$ is called \emph{big} if it is the complement of an analytic set of codimension at least $2$ in $Y$. Two line bundles on $Y$ which have isomorphic restrictions to a big open subset of $Y$ are in fact isomorphic over $Y$.

Let $\jepPubScript{F}$ be a subsheaf of $\jepPubScript{O}_Y(E)$. The complement of the biggest subset of $Y$ where $\jepPubScript{F}$ is the sheaf of sections of a subbundle $F$ of $E$ is called the \emph{singular locus} $\jepPubScript{S}(\jepPubScript{F})$ of $\jepPubScript{F}$. Equivalently, $\jepPubScript{S}(\jepPubScript{F})$ is the analytic subset of $Y$ where the quotient sheaf $\jepPubScript{O}_Y(E)/\jepPubScript{F}$ is not locally free. Recall that a subsheaf $\jepPubScript{F}$ of $\jepPubScript{O}_Y(E)$ is \emph{saturated} if $\jepPubScript{O}_Y(E)/\jepPubScript{F}$ is torsion free. The \emph{saturation} of $\jepPubScript{F}$ is the kernel of the morphism $\jepPubScript{O}_Y(E)\to(\jepPubScript{O}_Y(E)/\jepPubScript{F})_{\mathrm{tf}}$, where $(\jepPubScript{O}_Y(E)/\jepPubScript{F})_{\mathrm{tf}}$ denotes the quotient of $\jepPubScript{O}_Y(E)/\jepPubScript{F}$ by its torsion. If $\jepPubScript{F}$ is saturated, $\jepPubScript{S}(\jepPubScript{F})$ has codimension at least $2$ in $Y$, so that the \emph{regular locus} $Y\backslash\jepPubScript{S}(\jepPubScript{F})$ of $\jepPubScript{F}$ is a big open subset. Moreover a saturated subsheaf $\jepPubScript{F}$ is \emph{normal}, so that in particular for all big open subsets $\jepPubScript{U}$ of $Y$, the restriction $\jepPubScript{F}(Y)\to\jepPubScript{F}(\jepPubScript{U})$ is injective.

A subset of the foliated manifold $(Y,\jepPubScript{T})$ is said to be \emph{saturated under the foliation $\jepPubScript{T}$} if it is a union of leaves of $\jepPubScript{T}$. Obviously if $Z\subset Y$ is saturated under $\jepPubScript{T}$, then so is its complement $Y\backslash Z$. (Unfortunately, the word ``saturated'' is used here with two different meanings but this shouldn't cause any confusion.)

Here is the result (\cite[Prop.~2.2]{jep93KM17}).
\begin{proposition}
The harmonic Higgs bundle $(E,\theta)$ on the foliated Kähler manifold $(Y,\jepPubScript{T},\Omega)$ is weakly polystable along the leaves in the following sense:
\begin{enumerate}[label=(\arabic*)]
\item it is semistable along the leaves of $\jepPubScript{T}$: if $\jepPubScript{F}$ is a leafwise Higgs subsheaf of $(E,\theta)$, then $\deg_{\jepPubScript{T}}\jepPubScript{F}\leqslant0$.
\item if $\jepPubScript{F}$ is a saturated leafwise Higgs subsheaf of $(E,\theta)$ such that $\deg_{\jepPubScript{T}}\jepPubScript{F}=0$, then the singular locus $\jepPubScript{S}(\jepPubScript{F})$ of $\jepPubScript{F}$ is saturated under the foliation $\jepPubScript{T}$. Moreover, on $Y\backslash\jepPubScript{S}(\jepPubScript{F})$, and if $F$ denotes the holomorphic subbundle of $E$ such that $\jepPubScript{F}=\jepPubScript{O}_{Y\backslash\jepPubScript{S}(\jepPubScript{F})}(F)$ and $F^\perp$ its $C^\infty$ orthogonal complement w.r.t. the harmonic metric, then $\theta(F^\perp\otimes L)\subset F^\perp$. Finally, for each leaf $\jepPubLeaf{L}$ of $\jepPubScript{T}$ such that $\jepPubLeaf{L}\cap\jepPubScript{S}(\jepPubScript{F})=\varnothing$, $F^\perp\jepPubRestrict_{\jepPubLeaf{L}}$ is holomorphic on $\jepPubLeaf{L}$ and $(E,\theta)\jepPubRestrict_{\jepPubLeaf{L}}=(F,\theta\jepPubRestrict_F)\jepPubRestrict_{\jepPubLeaf{L}}\oplus(F^\perp,\theta\jepPubRestrict_{F^\perp})\jepPubRestrict_{\jepPubLeaf{L}}$ is a holomorphic orthogonal decomposition of Higgs bundles on $\jepPubLeaf{L}$.
\end{enumerate}
\end{proposition}
\subsection{The tautological foliation on the projectivized tangent bundle of a complex hyperbolic manifold}
An $n$-dimensional complex hyperbolic manifold $X$ is a quotient of the complex hyperbolic $n$-space $\mathbb H^n_{\mathbb C}=\operatorname{SU}(1,n)/\operatorname U(n)$ by a discrete torsion free subgroup $\Gamma$ of $\operatorname{SU}(1,n)$. It is a Hermitian locally symmetric space of rank 1. The complex hyperbolic space $\mathbb H^n_{\mathbb C}$ can be realized as the subset of negative lines in $\mathbb C^{n+1}$ for a Hermitian form $h$ of signature $(n,1)$ (recall our convention in Example~2.1), which is an open subset in the projective space $\mathbb{CP}^n$.

Intersections of lines of $\mathbb{CP}^n$ with $\mathbb H^n_{\mathbb C}$ are totally geodesic complex subspaces of $\mathbb H^n_{\mathbb C}$ isometric to the Poincaré disc. They are called \emph{complex geodesics}. The space $\jepPubScript{G}$ of complex geodesics is the homogeneous space $\operatorname{SU}(1,n)/\operatorname S(\operatorname U(1,1)\times\operatorname U(n-1))$. It is also a complex manifold, since it can be realized as the open subset of the Grassmannian manifold $\operatorname{Gr}_2(\mathbb C^{n+1})$ of $2$-planes in $\mathbb C^{n+1}$ consisting of planes on which the signature of $h$ is $(1,1)$.

The \emph{projectivized tangent bundle} $\mathbb PT_{\mathbb H^n_{\mathbb C}}$ of $\mathbb H^n_{\mathbb C}$ is the space of lines in the holomorphic tangent bundle $T_{\mathbb H^n_{\mathbb C}}$ of $\mathbb H^n_{\mathbb C}$. As a homogeneous space it can be identified with $\operatorname{SU}(1,n)/\operatorname S(\operatorname U(1)\times\operatorname U(1)\times\operatorname U(n-1))$. Again, its also a complex manifold (in fact a Kähler one) since it identifies with an open subset of the manifold $\mathrm F_{1,2}(\mathbb C^{n+1})$ of incomplete flags $(\ell\subset\Pi\subset\mathbb C^{n+1})$ with $\dim\ell=1$, $\dim\Pi=2$. A flag $(\ell\subset\Pi\subset\mathbb C^{n+1})$ belong to $\mathbb PT_{\mathbb H^n_{\mathbb C}}$ if the line $\ell$ is negative and the plane $\Pi$ has signature $(1,1)$ for the Hermitian form $h$.

The natural $\operatorname{SU}(1,n)$-equivariant fibration $\pi_{\jepPubScript{G}}:\mathbb PT_{\mathbb H^n_{\mathbb C}}\to\jepPubScript{G}$ which associates to a tangent line to $\mathbb H^n_{\mathbb C}$ the complex geodesic it defines is a disc bundle over $\jepPubScript{G}$.

By $\operatorname{SU}(1,n)$-equivariance, this fibration endows the projectivized tangent bundle $\mathbb PT_X=\Gamma\backslash\mathbb PT_{\mathbb H^n_{\mathbb C}}$ of $X=\Gamma\backslash\mathbb H^n_{\mathbb C}$ with a smooth holomorphic foliation $\jepPubScript{T}$ by complex curves, whose leaves are given by the tangent spaces of the (immersed) complex geodesics in $X$. This foliation is called the \emph{tautological foliation} of $\mathbb PT_X$ because the tangent line bundle $L$ to the leaves is naturally isomorphic to the tautological line bundle $\jepPubScript{O}_{\mathbb PT_X}(-1)$ on $\mathbb PT_X$.

The space $\jepPubScript{G}$ of complex geodesics of $\mathbb H^n_{\mathbb C}$ is a pseudo-Kähler manifold: it admits an $\operatorname{SU}(1,n)$-invariant Kähler form $\omega_{\jepPubScript{G}}$, which is closed, of type $(1,1)$, non degenerate, but not positive definite. The form $\omega_{\jepPubScript{G}}$ is moreover unique up to scaling. This form defines a transverse volume form $\Omega_{\jepPubScript{G}}$ for the tautological foliation $\jepPubScript{T}$ on $\mathbb PT_X$, and the associated notion of foliated degree $\deg_{\jepPubScript{T}}$ for sheaves on $\mathbb PT_X$ has the following fundamental property (after a suitable normalization of the involved forms) \cite[Prop.~3.1]{jep93KM17}.
\begin{proposition}
Assume that $X=\Gamma\backslash\mathbb H^n_{\mathbb C}$ is compact and let $\pi:\mathbb PT_X\to X$ be the projectivized tangent bundle of $X$. If $\jepPubScript{F}$ is a coherent $\jepPubScript{O}_X$-sheaf, then $\deg_{\jepPubScript{T}}(\pi^*\jepPubScript{F})=\deg_\omega\jepPubScript{F}$, where $\deg_\omega\jepPubScript{F}$ is the usual degree of $\jepPubScript{F}$ computed w.r.t. the Kähler form $\omega$ on $X$.
\end{proposition}

